% fig1_concept.tex  — Concept figure (DRAFT)
% Sim2Real multitask transfer for nanobody Tm + the three "what simulation works?" axes.
% Requires in preamble (see main.tex):
%   \usepackage{tikz}
%   \usetikzlibrary{positioning,arrows.meta,fit,backgrounds,shapes.geometric,calc}
%
% Include with:  % fig1_concept.tex  — Concept figure (DRAFT)
% Sim2Real multitask transfer for nanobody Tm + the three "what simulation works?" axes.
% Requires in preamble (see main.tex):
%   \usepackage{tikz}
%   \usetikzlibrary{positioning,arrows.meta,fit,backgrounds,shapes.geometric,calc}
%
% Include with:  % fig1_concept.tex  — Concept figure (DRAFT)
% Sim2Real multitask transfer for nanobody Tm + the three "what simulation works?" axes.
% Requires in preamble (see main.tex):
%   \usepackage{tikz}
%   \usetikzlibrary{positioning,arrows.meta,fit,backgrounds,shapes.geometric,calc}
%
% Include with:  % fig1_concept.tex  — Concept figure (DRAFT)
% Sim2Real multitask transfer for nanobody Tm + the three "what simulation works?" axes.
% Requires in preamble (see main.tex):
%   \usepackage{tikz}
%   \usetikzlibrary{positioning,arrows.meta,fit,backgrounds,shapes.geometric,calc}
%
% Include with:  \input{figures/fig1_concept}   inside a figure(*) environment.
\begin{tikzpicture}[
    font=\small,
    >={Stealth[length=2.2mm]},
    box/.style={draw, rounded corners=1.5pt, minimum height=7mm, minimum width=20mm,
                align=center, inner sep=3pt},
    enc/.style={box, fill=blue!12},
    shared/.style={box, fill=blue!22, minimum width=26mm},
    head/.style={box, fill=blue!30, minimum width=16mm},
    data/.style={box, fill=gray!10, minimum width=20mm},
    out/.style={box, fill=orange!25, minimum width=14mm},
    axis/.style={draw, rounded corners=2pt, fill=teal!8, align=left, inner sep=5pt,
                 minimum width=40mm, text width=40mm},
    flow/.style={->, thick, gray!70},
    lbl/.style={font=\footnotesize\itshape, gray!60!black},
]

% ---------- (a) Multitask transfer architecture ----------
\node[data] (seq) {VHH sequence\\\texttt{VQWATV...}};
\node[enc, right=9mm of seq] (esm) {ESM-2 encoder\\\scriptsize(frozen \emph{or} fine-tuned)};
\node[shared, right=9mm of esm] (mlp) {shared MLP\\\scriptsize(latent representation)};

% task heads (stacked)
\node[head, right=12mm of mlp, yshift=12mm] (htm)  {Tm head};
\node[head, right=12mm of mlp]             (hddg) {ddG head};
\node[head, right=12mm of mlp, yshift=-12mm](hmd)  {MD-feature\\head};

\node[out, right=8mm of htm]  (otm)  {$T_m$};
\node[out, right=8mm of hddg] (oddg) {$\Delta\Delta G$};
\node[out, right=8mm of hmd]  (omd)  {MD\,$Q$};

\draw[flow] (seq) -- (esm);
\draw[flow] (esm) -- (mlp);
\draw[flow] (mlp.east) -- (htm.west);
\draw[flow] (mlp.east) -- (hddg.west);
\draw[flow] (mlp.east) -- (hmd.west);
\draw[flow] (htm)  -- (otm);
\draw[flow] (hddg) -- (oddg);
\draw[flow] (hmd)  -- (omd);

% data-source annotations on the heads
\node[lbl, above=1mm of htm]  {experimental (scarce, $n\!\approx\!57$)};
\node[lbl, below=1mm of hmd]  {simulation (auxiliary, abundant)};

% brace grouping the auxiliary heads
\node[lbl, right=2mm of oddg, text width=22mm] {auxiliary\\transfer tasks};

% panel (a) label + title
\node[anchor=west, font=\bfseries] at ([yshift=10mm]seq.north west) {(a)\ \ Sim2Real multitask transfer learning};

% ---------- (b) The three design axes studied ----------
\node[axis, below=16mm of seq.south west, anchor=north west] (ax1)
    {\textbf{Source}\\physics $\Delta\Delta G$ (FEP/Rosetta)\\\emph{vs.} conventional MD $Q$};
\node[axis, right=6mm of ax1] (ax2)
    {\textbf{Temperature}\\$300\,$K \emph{vs.} $400\,$K\\(near vs.\ above $T_m$)};
\node[axis, right=6mm of ax2] (ax3)
    {\textbf{Feature}\\contact $Q$-value\\\emph{vs.} geometric (RMSF, SS, CDR3)};

\node[anchor=west, font=\bfseries] at ([yshift=7mm]ax1.north west)
    {(b)\ \ Which simulation data transfers? — three axes};

% ---------- (c) Evaluation: scaling + MRS ----------
\node[axis, fill=orange!8, below=10mm of ax2, anchor=north, text width=86mm, minimum width=86mm] (ev)
    {\textbf{(c)\ \ Evaluation:} data-scaling law $\mathrm{MAE}(n)=a\,n^{b}+c$ over \#simulation
     samples, compared against the \emph{experimental-data} scaling reference
     $\Rightarrow$ marginal rate of substitution (sim samples per 1 real $T_m$).};

\end{tikzpicture}
   inside a figure(*) environment.
\begin{tikzpicture}[
    font=\small,
    >={Stealth[length=2.2mm]},
    box/.style={draw, rounded corners=1.5pt, minimum height=7mm, minimum width=20mm,
                align=center, inner sep=3pt},
    enc/.style={box, fill=blue!12},
    shared/.style={box, fill=blue!22, minimum width=26mm},
    head/.style={box, fill=blue!30, minimum width=16mm},
    data/.style={box, fill=gray!10, minimum width=20mm},
    out/.style={box, fill=orange!25, minimum width=14mm},
    axis/.style={draw, rounded corners=2pt, fill=teal!8, align=left, inner sep=5pt,
                 minimum width=40mm, text width=40mm},
    flow/.style={->, thick, gray!70},
    lbl/.style={font=\footnotesize\itshape, gray!60!black},
]

% ---------- (a) Multitask transfer architecture ----------
\node[data] (seq) {VHH sequence\\\texttt{VQWATV...}};
\node[enc, right=9mm of seq] (esm) {ESM-2 encoder\\\scriptsize(frozen \emph{or} fine-tuned)};
\node[shared, right=9mm of esm] (mlp) {shared MLP\\\scriptsize(latent representation)};

% task heads (stacked)
\node[head, right=12mm of mlp, yshift=12mm] (htm)  {Tm head};
\node[head, right=12mm of mlp]             (hddg) {ddG head};
\node[head, right=12mm of mlp, yshift=-12mm](hmd)  {MD-feature\\head};

\node[out, right=8mm of htm]  (otm)  {$T_m$};
\node[out, right=8mm of hddg] (oddg) {$\Delta\Delta G$};
\node[out, right=8mm of hmd]  (omd)  {MD\,$Q$};

\draw[flow] (seq) -- (esm);
\draw[flow] (esm) -- (mlp);
\draw[flow] (mlp.east) -- (htm.west);
\draw[flow] (mlp.east) -- (hddg.west);
\draw[flow] (mlp.east) -- (hmd.west);
\draw[flow] (htm)  -- (otm);
\draw[flow] (hddg) -- (oddg);
\draw[flow] (hmd)  -- (omd);

% data-source annotations on the heads
\node[lbl, above=1mm of htm]  {experimental (scarce, $n\!\approx\!57$)};
\node[lbl, below=1mm of hmd]  {simulation (auxiliary, abundant)};

% brace grouping the auxiliary heads
\node[lbl, right=2mm of oddg, text width=22mm] {auxiliary\\transfer tasks};

% panel (a) label + title
\node[anchor=west, font=\bfseries] at ([yshift=10mm]seq.north west) {(a)\ \ Sim2Real multitask transfer learning};

% ---------- (b) The three design axes studied ----------
\node[axis, below=16mm of seq.south west, anchor=north west] (ax1)
    {\textbf{Source}\\physics $\Delta\Delta G$ (FEP/Rosetta)\\\emph{vs.} conventional MD $Q$};
\node[axis, right=6mm of ax1] (ax2)
    {\textbf{Temperature}\\$300\,$K \emph{vs.} $400\,$K\\(near vs.\ above $T_m$)};
\node[axis, right=6mm of ax2] (ax3)
    {\textbf{Feature}\\contact $Q$-value\\\emph{vs.} geometric (RMSF, SS, CDR3)};

\node[anchor=west, font=\bfseries] at ([yshift=7mm]ax1.north west)
    {(b)\ \ Which simulation data transfers? — three axes};

% ---------- (c) Evaluation: scaling + MRS ----------
\node[axis, fill=orange!8, below=10mm of ax2, anchor=north, text width=86mm, minimum width=86mm] (ev)
    {\textbf{(c)\ \ Evaluation:} data-scaling law $\mathrm{MAE}(n)=a\,n^{b}+c$ over \#simulation
     samples, compared against the \emph{experimental-data} scaling reference
     $\Rightarrow$ marginal rate of substitution (sim samples per 1 real $T_m$).};

\end{tikzpicture}
   inside a figure(*) environment.
\begin{tikzpicture}[
    font=\small,
    >={Stealth[length=2.2mm]},
    box/.style={draw, rounded corners=1.5pt, minimum height=7mm, minimum width=20mm,
                align=center, inner sep=3pt},
    enc/.style={box, fill=blue!12},
    shared/.style={box, fill=blue!22, minimum width=26mm},
    head/.style={box, fill=blue!30, minimum width=16mm},
    data/.style={box, fill=gray!10, minimum width=20mm},
    out/.style={box, fill=orange!25, minimum width=14mm},
    axis/.style={draw, rounded corners=2pt, fill=teal!8, align=left, inner sep=5pt,
                 minimum width=40mm, text width=40mm},
    flow/.style={->, thick, gray!70},
    lbl/.style={font=\footnotesize\itshape, gray!60!black},
]

% ---------- (a) Multitask transfer architecture ----------
\node[data] (seq) {VHH sequence\\\texttt{VQWATV...}};
\node[enc, right=9mm of seq] (esm) {ESM-2 encoder\\\scriptsize(frozen \emph{or} fine-tuned)};
\node[shared, right=9mm of esm] (mlp) {shared MLP\\\scriptsize(latent representation)};

% task heads (stacked)
\node[head, right=12mm of mlp, yshift=12mm] (htm)  {Tm head};
\node[head, right=12mm of mlp]             (hddg) {ddG head};
\node[head, right=12mm of mlp, yshift=-12mm](hmd)  {MD-feature\\head};

\node[out, right=8mm of htm]  (otm)  {$T_m$};
\node[out, right=8mm of hddg] (oddg) {$\Delta\Delta G$};
\node[out, right=8mm of hmd]  (omd)  {MD\,$Q$};

\draw[flow] (seq) -- (esm);
\draw[flow] (esm) -- (mlp);
\draw[flow] (mlp.east) -- (htm.west);
\draw[flow] (mlp.east) -- (hddg.west);
\draw[flow] (mlp.east) -- (hmd.west);
\draw[flow] (htm)  -- (otm);
\draw[flow] (hddg) -- (oddg);
\draw[flow] (hmd)  -- (omd);

% data-source annotations on the heads
\node[lbl, above=1mm of htm]  {experimental (scarce, $n\!\approx\!57$)};
\node[lbl, below=1mm of hmd]  {simulation (auxiliary, abundant)};

% brace grouping the auxiliary heads
\node[lbl, right=2mm of oddg, text width=22mm] {auxiliary\\transfer tasks};

% panel (a) label + title
\node[anchor=west, font=\bfseries] at ([yshift=10mm]seq.north west) {(a)\ \ Sim2Real multitask transfer learning};

% ---------- (b) The three design axes studied ----------
\node[axis, below=16mm of seq.south west, anchor=north west] (ax1)
    {\textbf{Source}\\physics $\Delta\Delta G$ (FEP/Rosetta)\\\emph{vs.} conventional MD $Q$};
\node[axis, right=6mm of ax1] (ax2)
    {\textbf{Temperature}\\$300\,$K \emph{vs.} $400\,$K\\(near vs.\ above $T_m$)};
\node[axis, right=6mm of ax2] (ax3)
    {\textbf{Feature}\\contact $Q$-value\\\emph{vs.} geometric (RMSF, SS, CDR3)};

\node[anchor=west, font=\bfseries] at ([yshift=7mm]ax1.north west)
    {(b)\ \ Which simulation data transfers? — three axes};

% ---------- (c) Evaluation: scaling + MRS ----------
\node[axis, fill=orange!8, below=10mm of ax2, anchor=north, text width=86mm, minimum width=86mm] (ev)
    {\textbf{(c)\ \ Evaluation:} data-scaling law $\mathrm{MAE}(n)=a\,n^{b}+c$ over \#simulation
     samples, compared against the \emph{experimental-data} scaling reference
     $\Rightarrow$ marginal rate of substitution (sim samples per 1 real $T_m$).};

\end{tikzpicture}
   inside a figure(*) environment.
\begin{tikzpicture}[
    font=\small,
    >={Stealth[length=2.2mm]},
    box/.style={draw, rounded corners=1.5pt, minimum height=7mm, minimum width=20mm,
                align=center, inner sep=3pt},
    enc/.style={box, fill=blue!12},
    shared/.style={box, fill=blue!22, minimum width=26mm},
    head/.style={box, fill=blue!30, minimum width=16mm},
    data/.style={box, fill=gray!10, minimum width=20mm},
    out/.style={box, fill=orange!25, minimum width=14mm},
    axis/.style={draw, rounded corners=2pt, fill=teal!8, align=left, inner sep=5pt,
                 minimum width=40mm, text width=40mm},
    flow/.style={->, thick, gray!70},
    lbl/.style={font=\footnotesize\itshape, gray!60!black},
]

% ---------- (a) Multitask transfer architecture ----------
\node[data] (seq) {VHH sequence\\\texttt{VQWATV...}};
\node[enc, right=9mm of seq] (esm) {ESM-2 encoder\\\scriptsize(frozen \emph{or} fine-tuned)};
\node[shared, right=9mm of esm] (mlp) {shared MLP\\\scriptsize(latent representation)};

% task heads (stacked)
\node[head, right=12mm of mlp, yshift=12mm] (htm)  {Tm head};
\node[head, right=12mm of mlp]             (hddg) {ddG head};
\node[head, right=12mm of mlp, yshift=-12mm](hmd)  {MD-feature\\head};

\node[out, right=8mm of htm]  (otm)  {$T_m$};
\node[out, right=8mm of hddg] (oddg) {$\Delta\Delta G$};
\node[out, right=8mm of hmd]  (omd)  {MD\,$Q$};

\draw[flow] (seq) -- (esm);
\draw[flow] (esm) -- (mlp);
\draw[flow] (mlp.east) -- (htm.west);
\draw[flow] (mlp.east) -- (hddg.west);
\draw[flow] (mlp.east) -- (hmd.west);
\draw[flow] (htm)  -- (otm);
\draw[flow] (hddg) -- (oddg);
\draw[flow] (hmd)  -- (omd);

% data-source annotations on the heads
\node[lbl, above=1mm of htm]  {experimental (scarce, $n\!\approx\!57$)};
\node[lbl, below=1mm of hmd]  {simulation (auxiliary, abundant)};

% brace grouping the auxiliary heads
\node[lbl, right=2mm of oddg, text width=22mm] {auxiliary\\transfer tasks};

% panel (a) label + title
\node[anchor=west, font=\bfseries] at ([yshift=10mm]seq.north west) {(a)\ \ Sim2Real multitask transfer learning};

% ---------- (b) The three design axes studied ----------
\node[axis, below=16mm of seq.south west, anchor=north west] (ax1)
    {\textbf{Source}\\physics $\Delta\Delta G$ (FEP/Rosetta)\\\emph{vs.} conventional MD $Q$};
\node[axis, right=6mm of ax1] (ax2)
    {\textbf{Temperature}\\$300\,$K \emph{vs.} $400\,$K\\(near vs.\ above $T_m$)};
\node[axis, right=6mm of ax2] (ax3)
    {\textbf{Feature}\\contact $Q$-value\\\emph{vs.} geometric (RMSF, SS, CDR3)};

\node[anchor=west, font=\bfseries] at ([yshift=7mm]ax1.north west)
    {(b)\ \ Which simulation data transfers? — three axes};

% ---------- (c) Evaluation: scaling + MRS ----------
\node[axis, fill=orange!8, below=10mm of ax2, anchor=north, text width=86mm, minimum width=86mm] (ev)
    {\textbf{(c)\ \ Evaluation:} data-scaling law $\mathrm{MAE}(n)=a\,n^{b}+c$ over \#simulation
     samples, compared against the \emph{experimental-data} scaling reference
     $\Rightarrow$ marginal rate of substitution (sim samples per 1 real $T_m$).};

\end{tikzpicture}
